\documentclass[man,english,donotrepeattitle]{apa7}
\usepackage[T1]{fontenc}
\usepackage[utf8]{inputenc}

% Font setup for PDFLaTeX - using Helvetica as Calibri alternative
\usepackage{helvet}
\renewcommand{\familydefault}{\sfdefault} % Set sans-serif as default

\usepackage{setspace}
\setstretch{1.3} % Set line spacing to 1.3
\usepackage{geometry}
\geometry{left=3cm, right=2.5cm, top=2.5cm, bottom=3cm}

% Page numbering setup
\usepackage{fancyhdr}
\usepackage{lastpage} % Required for X/Y page numbering
\pagestyle{fancy}
\fancyhf{}
\rfoot{\thepage/\pageref{LastPage}}

% Remove ALL headers and footers from title page
\fancypagestyle{plain}{%
  \fancyhf{}%
  \renewcommand{\headrulewidth}{0pt}%
  \renewcommand{\footrulewidth}{0pt}%
}

% Custom page style for title page that removes everything
\fancypagestyle{titlepage}{%
  \fancyhf{}%
  \renewcommand{\headrulewidth}{0pt}%
  \renewcommand{\footrulewidth}{0pt}%
}

% Enable section numbering (exclude title, abstract, and TOC)
\setcounter{secnumdepth}{3}

\usepackage{booktabs}
\usepackage{tabularx}
\usepackage{graphicx}
\usepackage{subcaption}
\usepackage{amsmath}
\usepackage[style=apa, backend=biber, natbib=true]{biblatex}
\addbibresource{references.bib}

\title{Illuminating Decisions: How Salient Reward-Paired Cues Modulate Learning and Choice Behavior}
\shorttitle{} % Remove short title to prevent header display
\author{Yannik Poth}
\affiliation{Rheinisch-Westfälische Technische Hochschule Aachen}
\course{Master of Science in Psychology} % Example, adjust as needed
\date{\today}

\begin{document}

\maketitle % Generates the title page
\thispagestyle{titlepage} % Use custom title page style with no numbering

\pagestyle{empty} % Prevent headers/footers on subsequent pages (abstract, TOC)

\newpage
\begin{center}
\textbf{Abstract}
\end{center}
\begin{abstract}
\noindent This study investigates how salient audiovisual feedback influences learning rates and choice behavior in a reinforcement learning framework. Using a restless two-armed bandit task, 44 participants made repeated choices between visual stimuli while salient reward-paired cues were systematically manipulated using a variable reinforcement schedule. A hierarchical Bayesian reinforcement learning model was employed to quantify learning dynamics, specifically examining whether salient feedback modulates the learning rate parameter ($\alpha_{\text{shift}}$). Additionally, win-stay probabilities were analyzed to assess simple behavioral strategies. Individual differences in impulsivity and sensation seeking were explored as potential moderators. The research aims to understand how sensory salience affects learning processes and decision-making, with implications for understanding vulnerability to environments designed to maximize engagement, such as electronic gambling machines. Results provide insights into the mechanisms by which environmental design features can influence adaptive learning processes and potentially contribute to maladaptive choice patterns.
\end{abstract}

% Table of Contents
\newpage
\tableofcontents
\newpage

\pagestyle{fancy} % Restore fancy page style for the main document body

% Introduction with proper heading
\section{Introduction}
In modern society, we are constantly surrounded by environments designed to capture our attention through potent sensory enticements—such as bright lights, vivid colors, and distinctive sounds. From smartphone notifications to dynamic digital advertisements, these attention-grabbing stimuli have become ubiquitous in our daily lives. Perhaps nowhere is this phenomenon more deliberately engineered than in gambling environments, particularly in electronic gambling machines (EGMs). When a player wins on a digital slot machine, they are greeted not merely with monetary reward, but with an orchestrated sensory celebration—flashing lights, triumphant sounds, and animated displays—creating a multisensory experience far more salient than the simple fact of winning money. This deliberate design strategy is not merely aesthetic; research indicates that these audiovisual features may influence gambling persistence—especially for higher‑risk gamblers—with players demonstrating longer play sessions on machines that provide enhanced sensory feedback.

This strategic pairing of rewards with intensified sensory feedback raises a fundamental question about human cognition: How do salient reward-paired cues (SRC) influence the way we learn from outcomes and make subsequent decisions? The answer to this question has implications extending far beyond gambling contexts—it touches upon basic mechanisms of learning, decision-making, and ultimately behaviors that can range from adaptive to potentially harmful. From a clinical perspective, understanding these mechanisms may provide insights into vulnerability factors for conditions like gambling disorder, which affects approximately 1.2\% of the population and is associated with significant psychological, social, and financial harm \citep{who2024}. Therefore, this thesis investigates how such SRCs, inspired by EGM design, impact learning dynamics and choice behavior in a controlled setting.

\subsection{Reinforcement Learning as a Framework for Understanding Decision-Making}

Reinforcement learning (RL) provides a powerful and computationally precise framework for understanding how humans and other animals learn from the consequences of their actions. At its core, RL posits that agents adjust their expectations and behaviours based on prediction errors—the discrepancy between expected and received outcomes \citep{sutton2020}. When we experience an unexpected reward, a positive prediction error signal drives us to increase the value we assign to the action that led to that reward, making us more likely to repeat it in the future. Conversely, unexpected negative outcomes lead us to devalue certain choices and avoid them subsequently.

A critical parameter in reinforcement learning models is the learning rate (alpha, $\alpha$), which determines how quickly value estimates are updated in response to new information. Formally, the value update can be expressed as: 
\begin{equation*}
V_{\text{new}} = V_{\text{old}} + \alpha \times (R - V_{\text{old}})
\end{equation*}
where $V$ represents the value estimate before ($V_{\text{old}}$) and after ($V_{\text{new}}$) receiving the outcome, $R$ is the outcome received, $\alpha$ represents the learning rate, typically constrained between 0 and 1 \citep{sutton2020}. The term $(R-V_{\text{old}})$ quantifies the prediction error, representing the difference between the actual reward obtained and the value that was expected. Higher learning rates result in more rapid updating based on recent outcomes, while lower rates produce more stable value estimates that change gradually over time. Importantly, research suggests that learning rates are not fixed but can be dynamically modulated by contextual factors, emotional states, and attentional processes \citep{lepelley2016, pike2022, wu2017}. This adaptive flexibility in learning rates may be advantageous in volatile environments but could potentially be exploited by artificial environments designed to maximize engagement.

Beyond these formal value-updating mechanisms, decision-making can sometimes be guided by simpler heuristics. One such fundamental strategy is ``win-stay, lose-shift'' (WSLS). This heuristic dictates a straightforward policy: if a chosen action leads to a positive outcome (a win), repeat that action on the next opportunity; if it leads to a negative outcome (a loss), switch to an alternative action \citep{robbins1952, nowak1993}. While computationally less demanding than maintaining and updating precise value estimates, WSLS can be surprisingly effective, particularly in environments where outcomes are somewhat stable. In the context of the current study, examining the win-stay component of this strategy provides a direct behavioral index of how recent rewards influence subsequent choices. It is therefore pertinent to investigate whether factors like feedback salience modulate not only the computationally derived learning rate ($\alpha$) but also this more basic tendency to repeat rewarded actions.

Traditional RL models often focus on the valence (positive/negative) and magnitude of outcomes, potentially overlooking how the sensory presentation of these outcomes—their salience—might directly influence the learning process. Salience refers to the degree to which a stimulus stands out and captures attention, which can be due to its sensory intensity (e.g., brightness, loudness), novelty, or motivational significance \citep{berridge1998, itti2001}. Empirical studies have demonstrated that manipulating the visual salience of stimuli can significantly bias attention allocation and subsequent decision-making \citep{towal2013}, and that salient perceptual features can override value-based decision processes under certain conditions \citep{tsetsos2012}.

\subsection{Electronic Gambling Machines and Psychological Salience}

The gambling industry appears to intuitively understand the power of salient feedback, as evidenced by the sophisticated sensory design of modern EGMs. These machines create an immersive environment where wins—and notably, ``losses disguised as wins'' (LDWs)—are celebrated with disproportionate audiovisual feedback. LDWs occur when outcomes that are net losses (smaller than the initial wager) are nevertheless accompanied by celebratory sensory stimuli typically associated with wins, confusing players about their actual financial outcomes and contributing to an inflated sense of winning frequency \citep{dixon2010}. This approach exploits cognitive mechanisms whereby salient feedback enhances the effectiveness of rewards in driving learning and decision-making.

Experimental evidence directly supports this design strategy. Research by \citet{mereu2010} found that visually enhancing win outcomes significantly increased participants' illusion of control—their perception of influence over random outcomes—while similarly enhanced loss feedback had minimal effect. Indeed, research suggests that EGMs with enhanced audiovisual feedback—and especially LDWs—promote persistent gambling even in the face of net losses presumably by increasing play desirability and misperceptions of wins, though effects vary by individual \citep{dixon2014, graydon2018, myles2023}. This suggests that SRCs may not only intensify the experience of genuine rewards but potentially distort the very categorization of outcomes in decision-making processes.

This strategic deployment of salience has intensified with the technological evolution of EGMs, which now incorporate highly sophisticated audiovisual features synchronized with game events \citep{schull2014}. Comparative studies demonstrate that machines with enhanced sensory features, such as richer sounds, are often preferred by players and generate higher revenue, even when the underlying probabilities and payout structures are identical \citep{dixon2014, graydon2018}. In a controlled laboratory experiment, \citet{cherkasova2018} demonstrated that SRCs directly promoted riskier choices in a gambling task, independent of the options' expected values, providing evidence for how such features might influence decision-making. In practical terms, participants essentially chose to play against the mathematical odds simply due to the enhanced sensory experience of winning.

Furthermore, physiological measurements indicate that SRCs elicit stronger autonomic responses, including increased skin conductance and heart rate deceleration—markers of heightened attention and emotional arousal \citep{dixon2014, detez2019}. Understanding the precise cognitive and neural mechanisms driving these effects is therefore crucial.

\subsection{Neurocognitive Mechanisms of Salience and Reward Processing}

Exploring the neurocognitive underpinnings reveals how salience might interact with reward processing. From a cognitive neuroscience perspective, this phenomenon involves multiple brain systems operating in parallel. Research indicates that the brain processes outcome information through distinct yet interacting pathways—an early salience detection system involving regions such as the anterior insula and dorsal anterior cingulate cortex, which are critical for identifying behaviorally relevant stimuli \citep{seeley2007}, followed by value computation in the ventral striatum and ventromedial prefrontal cortex \citep{bartra2013}.

Human fMRI work shows that when monetary rewards are rendered behaviorally salient, activity in the ventral striatum is amplified \citep{zink2004}. At the cellular level, a subset of mid‑brain dopaminergic neurons displays stronger phasic bursts when rewards (or other outcomes) are preceded by novel or attention‑capturing cues, indicating that salience can boost dopaminergic signaling beyond classical value coding, potentially enhancing the effective learning rate beyond what reward magnitude alone would predict \citep{bromberg2010}. Consistent with this, fibre‑photometry recordings reveal larger nucleus‑accumbens dopamine transients as the physical intensity—and hence perceived salience—of appetitive or aversive stimuli increases \citep{kutlu2021}. Collectively, these findings point towards plausible neural mechanisms for the proposed $\alpha$ shift: salient feedback may amplify dopaminergic prediction error signals, leading to a larger update step and thus a higher effective learning rate specifically for saliently signaled outcomes.

Beyond direct modulation of reward signals, SRCs also influence attention and arousal. In the aforementioned study by \citet{cherkasova2018}, eye-tracking revealed that participants spent less time examining probability information when salient win cues were present, suggesting a shift in attentional priorities away from deliberative value assessment. Concurrently, pupillometry indicated greater pupil dilation following wins in the cued condition. This pattern, linked to heightened locus coeruleus-norepinephrine activity, suggests increased physiological arousal elicited by the salient rewards. Together, these findings indicate that sensory salience can shift both what information is prioritized and how strongly rewarding outcomes shape subsequent choices—offering a plausible mechanism through which salience may modulate learning.

\subsection{Individual Differences in Susceptibility to Salient Cues}

An important consideration is that the impact of SRCs likely varies across individuals. Personality traits such as sensation seeking—characterized by the desire for novel, varied, and intense sensory experiences—and impulsivity—the tendency toward rapid, unplanned actions with diminished regard for negative consequences—have been linked to both gambling behavior and heightened neural responses to reward‑predictive cues \citep{hawes2017, verdejo2008, zuckerman1994}. These traits may moderate the effect of salient feedback on learning rates and decision strategies, potentially explaining why some individuals are more vulnerable to environments rich in such cues.

Indeed, higher sensation seeking is associated with stronger psychophysiological and neural responses to reward cues \citep{hawes2017} and potentially biased learning, favoring positive over negative outcomes \citep{xu2019}. Trait impulsivity has also been linked to altered RL profiles, sometimes manifesting as reduced learning from negative feedback compared to rewards \citep{pauli2025, zou2022}. Clinical data also confirm that higher impulsivity co‑occurs with greater gambling severity \citep{blaszczynski1997, slutske2005}. Therefore, these individual differences could interact with environmental factors like feedback salience, potentially amplifying or dampening its effect on learning and contributing to varied susceptibility to maladaptive decision-making in stimulating environments.

\subsection{Research Gaps and Current Study}

While the research reviewed above highlights the potent influence of salient sensory feedback, particularly in contexts like EGM play, several gaps remain in understanding the precise mechanisms through which these cues modulate learning and decision-making. Specifically, further research is needed to quantify these effects within a formal learning framework and to isolate the impact of salience itself.

First, although behavioral consequences like increased gambling persistence are well-documented \citep{dixon2014, graydon2019}, fewer studies have employed computational modeling, such as reinforcement learning, to formally quantify how salient feedback specifically alters underlying learning parameters like the learning rate ($\alpha$). Understanding whether salience directly modifies the weight given to prediction errors is crucial for mechanistic insight. Furthermore, existing research has typically compared machines with different levels of overall sensory stimulation rather than systematically manipulating the presence of salient feedback while holding reward values constant. This methodological limitation has made it difficult to isolate the specific effect of sensory salience on learning processes.

Second, the choice of task is critical. While effects have been observed in various paradigms \citep{cherkasova2018, mereu2010}, tasks specifically designed to probe trial-by-trial adjustments based on feedback, such as multi-armed bandit tasks, are particularly well-suited for quantifying learning rate dynamics using RL models. The use of a restless bandit task, where reward probabilities change over time, further necessitates continuous learning and adaptation, making it sensitive to factors modulating learning speed.

Third, while the general impact of salience on attention and choice is studied, the specific way SRCs—a hallmark of EGM design—influence reinforcement learning parameters and simple behavioral strategies like win-stay behavior requires more targeted investigation in controlled experimental settings.

Finally, although individual differences in traits like sensation seeking and impulsivity are known moderators of reward sensitivity and gambling behavior \citep{hawes2017, pauli2025}, how these traits specifically relate to the magnitude of learning modulation induced by salient feedback (i.e., the $\alpha$ shift) remains an open area for investigation.

The current study is designed to address these gaps using a controlled experimental approach. Specifically, a restless two-armed bandit task is employed where participants make choices under conditions of changing reward probabilities. Crucially, within this task, the presence of salient audiovisual feedback accompanying win outcomes is systematically manipulated using a variable ratio schedule, while the reward outcomes themselves remain independent of the feedback type. This within-subject design allows for direct comparison of behavior and learning parameters between trials with salient versus non-salient feedback for identical win outcomes, thereby isolating the effect of sensory salience.

The primary objective is to investigate how salient, win-paired audiovisual cues modulate learning rates and choice strategies within an RL framework. To quantify these effects, a hierarchical Bayesian reinforcement learning model will be fitted to the trial-by-trial choice data. It is hypothesized that:

\begin{enumerate}
\item The presence of salient feedback following wins will enhance learning from those outcomes, manifesting as a positive $\alpha_{\text{shift}}$ parameter within the RL model. This parameter explicitly captures the difference in effective learning rates between salient and non-salient win trials.

\item Salient feedback following a win will increase the likelihood of repeating the rewarded choice on the subsequent trial, resulting in higher win-stay probability compared to wins followed by non-salient feedback.
\end{enumerate}

Furthermore, the study will explore potential associations between individual differences in sensation seeking and impulsivity and the magnitude of the observed effects of salient feedback on learning ($\alpha_{\text{shift}}$) and choice behavior (win-stay probability).

By integrating computational modeling with a precisely controlled experimental design targeting dynamic learning, this research aims to provide a clearer understanding of how salient environmental cues, inspired by real-world applications like EGMs, can shape adaptive learning processes and potentially contribute to maladaptive choice patterns. 

\section{Method}
% Method section content

\subsection{Participants}

The final analyzed sample consisted of 44 individuals. This group comprised 22 participants identifying as female, 21 as male, and one as non-binary. Participant ages ranged from 18 to 54 years ($M = 24.86$, $SD = 7.02$). Nearly half of the participants (47.73\%) were enrolled in psychology programs at the University of Cologne.

Inclusion criteria for this study specified an age range of 18–60 years and no self-reported history of current or past neurological or psychiatric disorders. The exclusion of individuals with neurological or psychiatric conditions was implemented to ensure sample homogeneity and the validity of findings related to reinforcement learning, as these conditions can significantly alter underlying neural and cognitive processes \citep{pike2022}; this criterion also served an ethical purpose by avoiding potential additional stress for vulnerable individuals. The age range was established to minimize potential confounds from age-related changes in dopaminergic systems critical for reinforcement learning \citep{bolenz2019}. Individuals outside these criteria, or those who self-reported other conditions that might substantially interfere with task performance, were excluded. All 44 recruited participants met these criteria and completed the study.

The target sample size was determined a priori using G*Power (Version 3.1; \citealt{faul2007}). An anticipated moderate effect size of Cohen's $d = 0.4$ was used for this calculation. This effect size was selected based on a review of existing literature investigating the influence of salient sensory cues on decision-making and learning processes. While direct research on the precise impact of salient feedback on learning rate parameters is still developing, studies on related phenomena—such as the promotion of riskier choices by SRCs \citep{cherkasova2018, cherkasova2024} or the modulation of decision strategies by salient information \citep{liu2018, mereu2010}—report moderate effect sizes, typically ranging from Cohen's $d \approx 0.5$ to 0.7. A slightly more conservative estimate of Cohen's $d = 0.4$ was chosen for the present study to account for potential individual differences in susceptibility to salient cues and the specific focus on learning rate modulation, rather than broader choice behavior. To detect an effect of this magnitude for the primary directional hypothesis regarding learning rates (Hypothesis 1.1) with $\alpha = .05$ and 80\% power ($1 - \beta = .80$) for a one-tailed dependent means $t$-test, a minimum of 41 participants was required. This was rounded to 42 to support gender-balanced recruitment, with non-binary individuals also welcomed. Although the primary analyses utilize Bayesian hierarchical modeling, this conventional power analysis informed the sample size target. The final sample of 44 participants exceeded this minimum.

Recruitment utilized multiple channels. Primarily, participants were sourced from the University of Cologne's subject pool, with study advertisements disseminated through the study portal of the Faculty of Human Sciences and institutional email distribution lists. These efforts were supplemented by physical posters displayed on campus, outreach within student WhatsApp groups, and word-of-mouth referrals.

Participation in the study lasted approximately 40 minutes. As compensation, individuals selected either course credit, equivalent to 0.75 participant hours for University of Cologne students, or a monetary payment of €8, calculated at a rate of €12 per hour. Furthermore, all participants were eligible for a performance-based cash bonus of up to €3, contingent on points earned during the task and awarded independently of their chosen base compensation.

\subsection{Materials}

\subsubsection{Task Apparatus}

The experiment was administered on a desktop computer equipped with a 24-inch LCD monitor displaying stimuli at a resolution of 1920×1080 pixels. Participants wore Sennheiser HD 280 Pro closed-back headphones for the delivery of auditory stimuli throughout the task.

The experimental task was developed and executed within a Python 3.10.0 (64-bit) environment. Precise visual stimulus presentation, timing control, and response collection were managed using PsychoPy (v2024.2.4; \citealt{peirce2019}). Auditory feedback, including continuous background music and distinct salient feedback sounds, was handled via Pygame (v2.6.1).

\subsubsection{Task Stimuli}

In each trial of the main task, participants chose between two distinct fractal images. These images were generated using DALL-E 3 \citep{openai2024} to serve as unique choice options. The stimuli shared several core visual characteristics, including intricate, high-contrast symmetrical designs, a luminous quality that created a sense of depth, and a clear central emphasis. Both images were consistently displayed within identical rounded square frames featuring metallic gradient borders. While similar in fundamental visual properties and presentation, the two images primarily differed in their dominant color and overall pattern, with one featuring a red, radial or starburst design, and the other a blue, geometric grid conveying an architectural appearance. For the practice trials, two similarly constructed and framed fractal images were utilized, one colored green and the other purple.

For all choices, an immediate visual confirmation was provided by a simple border appearing around the selected stimulus paired with immediate text feedback informing about win or loss. On designated win trials, this standard feedback was augmented by a salient visual display. A three-second animation of the chosen stimulus, inspired by electronic gambling machine visuals and created by the author using Apple Motion (v5.7), featured a pulsating, blooming effect with dynamic lighting and contrast shifts designed to accentuate the fractal's structure, enhance depth and movement, and thereby maximize visual salience.

The auditory environment was characterized by continuous, ambient, and spacey background music track played throughout the task. The track has a duration of 105 seconds that looped seamlessly. It featured a sustained, modulated dynamic profile intended to be consistent yet non-intrusive. Distinctly, salient visual animations were paired with a three-second, high-amplitude, explosive sound effect, designed as an exaggerated, attention-capturing auditory cue. Both the background music and the salient sound effect were composed by the author using Apple Logic Pro (v10.7.7).

\subsubsection{Questionnaires}

Following the experimental task, participants completed a series of questionnaires presented digitally via a custom Python program. These included standardized measures of personality traits and a study-specific evaluation questionnaire.

To assess impulsivity, the short German version of the Barratt Impulsiveness Scale (BIS-15; \citealt{meule2011}) was administered. This 15-item self-report measure evaluates impulsivity across three subdimensions: non-planning, motor, and attentional impulsivity. Participants respond on a 4-point Likert scale, with higher total scores indicating greater levels of impulsivity. The German BIS-15 has demonstrated good reliability, with a Cronbach's alpha of .81 for the total scale, and has established validity in German-speaking samples \citep{meule2011}.

Sensation seeking was measured using an 8-item German short version of the Sensation Seeking Scale (SSS; \citealt{grimm2015}), adapted from Zuckerman's original work \citep{zuckerman1979}. This instrument assesses four facets: Thrill and Adventure Seeking, Experience Seeking, Disinhibition, and Boredom Susceptibility. For each item, participants choose between two statements. One item within the Experience Seeking subdimension was adapted by the author to replace culturally insensitive language with contemporary, non-stigmatizing terms, while aiming to preserve the item's original intent of assessing openness to unconventional groups. It is noted that this adapted short version lacks formal psychometric validation data.

Finally, participants completed a custom-designed evaluation questionnaire to provide feedback on their experience during the study. This brief questionnaire comprised nine items, employing a mix of open-ended and closed-ended questions. The open-ended items aimed to explore participants' understanding of the study's overall goal and their subjective perceptions of the task, particularly regarding the audiovisual feedback. Closed-ended items assessed awareness of the salient feedback manipulations and their perceived influence on decision-making and motivation. The responses to this questionnaire were intended for qualitative exploration and were not part of the primary quantitative analyses.

\subsection{Procedure}

Upon arrival, participants provided written informed consent and were seated in a quiet, isolated testing room with headphones. A custom Python script administered the entire experimental session. On-screen instructions introduced the restless two-armed bandit task. These explained that the objective was to accumulate points by selecting between two visual options using arrow keys, with immediate feedback following each choice. The instructions emphasized that reward probabilities would fluctuate throughout the task, requiring continuous adaptation of choice strategy. They clarified that the visual symbol chosen was relevant rather than its screen position. Participants were also informed, that some wins would feature ``special effects''—salient audiovisual feedback—and that performance could earn a bonus of up to €3.00.

A 15-trial practice phase followed to familiarize individuals with the task mechanics, including both non-salient and salient feedback types. To mitigate novelty effects, the salient audiovisual feedback during practice followed the same reinforcement schedule as the main task. Different stimuli were used for practice trials to prevent stimulus-specific learning from transferring.

After a brief on-screen reminder of task goals, the main experimental phase commenced with 200 trials, a number selected to enable robust analysis of learning dynamics under different feedback conditions. Throughout both phases, an experimenter remained available to address any technical issues. Upon task completion, digital questionnaires were administered in fixed order: the Barratt Impulsiveness Scale (BIS-15), the Sensation Seeking Scale (SSS), and an experiment-specific evaluation questionnaire. This sequence was designed to assess stable personality traits before exploring specific experimental experiences. The session concluded with a brief debriefing about the study's purpose—investigating how salient audiovisual feedback influences learning and decision-making. Each participant received their predetermined compensation after being thanked for their involvement.

\subsection{Task Design}

Participants engaged in a restless two-armed bandit task, requiring them to make repeated choices between two visual stimuli to maximize their accumulated points.

\subsubsection{Basic Structure and Trial Events}

On each trial, two distinct images (as described in 3.2.2 Task Stimuli) were presented concurrently on the screen. The on-screen positions (left or right) of these two stimuli were randomized on each trial to ensure that participants learned to associate rewards with the identity of the chosen stimulus rather than its spatial location. Choices were made using the left and right arrow keys. Internally, one specific stimulus was consistently mapped to a choice code of `0' and the other to `1' throughout the main task for each participant.

The sequence of events within a trial was precisely timed. The two choice stimuli were displayed until the participant responded or a 5-second response window expired. If no response was made within this window, the trial was marked as missed and no reward was delivered. Following a valid choice, feedback was presented immediately and remained on screen for 3 seconds. After the feedback offset, a central fixation cross (`+') appeared for a variable inter-trial interval (ITI), drawn from a normal distribution with a mean of 1 second and a standard deviation of 0.5 seconds (minimum ITI: 0.5 seconds), before the next trial commenced.

\subsubsection{Reward Structure and ``Restless'' Probabilities}

A choice on any given trial resulted in a binary outcome: either a reward (win) or no reward (loss). The binary structure was deliberately selected to provide clear win/no-win signals, thereby simplifying the learning problem. This methodological approach allows the investigation to focus specifically on how salience modulates the processing of discrete outcomes, rather than examining responses to varying reward magnitudes.

The experimental design employed a dynamic probabilistic structure wherein reward probabilities for both options varied independently and continuously, requiring participants to constantly update their value estimations. These underlying reward probabilities were generated using a decaying Gaussian random walk model, consistent with established methodologies in the field \citep{daw2006}. Initial reward probabilities for each option were drawn from a Gaussian distribution ($\mu_{t_0} = 0.5$; $\sigma_o = 0.04$). Subsequently, these probabilities evolved according to a formula $\mu_{i,t+1} = \lambda\mu_{i,t} + (1-\lambda)\theta + \nu_t$, where $\mu_{i,t}$ represents the reward probability for option $i$ at trial $t$, $\lambda = 0.9836$ is the decay parameter causing gradual reversion toward the decay center $\theta = 0.5$, and $\nu_t$ is Gaussian diffusion noise ($\mu = 0$, $\sigma_d = 0.065$). To increase environmental volatility, the diffusion noise terms between the two options were partially negatively correlated (corr = 0.35), with all probabilities constrained within the [0, 1] range. To prevent boundary stagnation and maintain participant engagement, both options' probabilities were periodically reset to 0.5 at irregular intervals (approximately every 30 trials on average, with temporal variability). The actual reward delivery on any given trial was determined stochastically by comparing the selected option's true underlying reward probability against a randomly generated number.

\subsubsection{Feedback Conditions}

When a choice resulted in a reward, the feedback provided was either non-salient or salient. The type of feedback delivered upon a win was determined by a variable reinforcement pattern. This pattern was generated by repeatedly shuffling a base binary sequence (e.g., [0, 0, 0, 1], where 1 signified salient feedback eligibility and 0 non-salient) and then concatenating these shuffled blocks. This procedure ensured that salient feedback occurred, on average, for one out of every four wins, but in an unpredictable, variable sequence across rewarded trials. The visual and auditory details of non-salient and salient feedback are described in the Materials section (3.2.2 Task Stimuli). Regardless of whether feedback was non-salient or salient, winning choices were always accompanied by on-screen text indicating ``+100 Punkte'' (German for ``+100 points''). If a choice resulted in a loss, the visual feedback was consistently non-salient, involving a static border around the chosen stimulus and on-screen text indicating ``+0 Punkte'' (German for ``+0 points''); in such instances, only the continuous background music was audible, with no specific loss-related sound effect. Missed trials, occurring when participants failed to respond within the 5-second window, triggered a ``Zu spät'' message (German for ``too late'') and were distinctly coded.

\subsubsection{Data Collection}

Comprehensive data were recorded meticulously for each trial. This included participant ID, experimental phase (practice or main), trial number, the participant's choice (coded 0 or 1, or None if missed), reaction time, reward outcome (1 for win, 0 for loss), feedback condition code (0 for non-salient, 1 for salient, 2 for missed trial), and the true underlying reward probabilities for both options on that trial. This trial-by-trial data was saved in real-time to a comma-separated values (CSV) file, uniquely named for each participant.

\subsection{Ethical Considerations}

The study protocol, encompassing all materials and procedures, received approval from the Ethics Committee of the Faculty of Human Sciences, University of Cologne (Ethikkommission der Humanwissenschaftliche Fakultät, Universität zu Köln; Approval No. JPHF0271, granted on December 11, 2024). All participants provided written informed consent before participation. Data were collected and stored anonymously, using unique participant IDs to ensure confidentiality and compliance with General Data Protection Regulation (GDPR), as detailed in the study's preregistration \citep{poth2025}.

\subsection{Statistical Approach}

Data preprocessing, visualization, and statistical analyses were conducted using R (Version 4.4.1; \citealt{rcore2024}). The hierarchical Bayesian models were specified and fitted using Stan (Version 2.32.2; \citealt{stan2025}), interfaced from R with the \texttt{rstan} package (Version 2.32.6; \citealt{standev2024}).

\subsubsection{Data Preprocessing}

All primary analyses were focused on data from the 200 main experimental trials per participant, excluding the 15 practice trials. Trials where participants failed to respond within the 5-second window were treated as missing data. These missed trials were coded with sentinel values (-9) and were subsequently ignored in the likelihood calculation of the reinforcement learning models, ensuring they did not contribute to parameter estimation. No other systematic data cleaning was applied before the primary model fitting.

\subsubsection{Model-Agnostic Behavioral Analysis}

To provide model-agnostic insights into choice behavior, a pre-registered analysis was conducted alongside two exploratory extensions.

The primary analysis directly tested Hypothesis 1.2 through win-stay probability assessment. Win-stay probability—the conditional probability of selecting the same option on trial $t+1$ following its reward on trial $t$—was calculated for each participant separately for trials following non-salient wins versus salient wins. A paired-samples $t$-test was specified to compare these conditions, with the Wilcoxon signed-rank test as a non-parametric alternative if normality assumptions were violated.

As exploratory extensions not specified in the pre-registration, two logistic generalized linear mixed-effects models (GLMMs) were fitted to leverage trial-by-trial data. The first GLMM extended the win-stay analysis by modeling stay probability across all three previous outcomes: loss, non-salient win, and salient win. This provided model-based estimates of both win-stay and lose-stay probabilities. The second GLMM examined choice optimality, using a binary indicator of whether choices were directed toward the stimulus with the current higher true reward probability.

Both GLMMs included the outcome of the previous trial and trial number as fixed effects, with random intercepts for each participant to account for individual differences. This approach utilized the complete dataset while accounting for within-subject dependencies.

\subsubsection{Hierarchical Bayesian Modeling}

In addition to the analysis of model-agnostic performance indices, the main focus was on computational modeling using a hierarchical Bayesian approach. A range of competing computational accounts of behavior was compared to identify the most parsimonious model explaining the data.

The models were constructed from a common set of components: a learning rule to update stimulus values and a choice rule to generate action probabilities. They differ in their complexity regarding the influence of feedback salience and choice history.

The \textit{learning rule} is based on the Rescorla-Wagner model \citep{rescorla1972}, where the expected value (Q-value) for a chosen stimulus $i$ on trial $t$ is updated based on a prediction error:
\begin{equation*}
    Q_{t+1}(i) = Q_t(i) + \alpha_{\text{eff}} \cdot \delta_t
\end{equation*}
The prediction error, $\delta_t = r_t - Q_t(i)$, is the discrepancy between the reward received ($r_t \in \{0, 1\}$) and the expected value. The effective learning rate, $\alpha_{\text{eff}}$, is modulated by feedback salience. To constrain learning rates to the [0, 1] interval, raw parameters are transformed using the standard normal cumulative distribution function ($\Phi$):
\begin{equation*}
    \alpha_{\text{eff}} = 
    \begin{cases} 
        \Phi(\alpha_{\text{raw}} + \alpha_{\text{shift\_raw}}) & \text{if feedback is salient} \\
        \Phi(\alpha_{\text{raw}}) & \text{if feedback is non-salient}
    \end{cases}
\end{equation*}
Initial Q-values for both stimuli were set to 0.5, representing no initial preference.

The \textit{choice rule} uses a softmax function to transform Q-values into choice probabilities. The probability of selecting stimulus $i$ on trial $t$ is given by:
\begin{equation*}
    P(i_t = i) = \frac{\exp(\beta \cdot Q_t(i))}{\sum_{j=1}^{2} \exp(\beta \cdot Q_t(j))}
\end{equation*}
where the inverse temperature parameter, $\beta$, controls choice stochasticity. It is transformed from an unbounded raw parameter to the range [0, 10] via $\beta = \Phi(\beta_{\text{raw}}) \cdot 10$.

To account for choice repetition tendencies, an extension to the choice rule includes a \textit{perseveration bonus}, $\kappa$. This parameter modifies the choice logits as follows:
\begin{equation*}
    P(i_t = i) = \frac{\exp(\beta \cdot Q_t(i) + \kappa \cdot I(i = c_{t-1}))}{\sum_{j=1}^{2} \exp(\beta \cdot Q_t(j) + \kappa \cdot I(j = c_{t-1}))}
\end{equation*}
where $I(x)$ is an indicator function that equals 1 if condition $x$ is true, and 0 otherwise. A positive $\kappa$ indicates a tendency to perseverate, while a negative $\kappa$ indicates a switching bias.

Three competing models were specified by combining these components:
\begin{enumerate}
    \item \textbf{M1 (Baseline):} A standard Rescorla-Wagner model, using the basic learning rule (i.e., $\alpha_{\text{shift\_raw}}$ is fixed to 0) and the basic softmax choice rule. This model has two free parameters: $\alpha$ and $\beta$.
    \item \textbf{M2 (Salience-Shift):} The primary model of interest, which uses the salience-modulated learning rule and the basic softmax choice rule. It has three free parameters: $\alpha$, $\alpha_{\text{shift}}$, and $\beta$.
    \item \textbf{M3 (Salience-Shift + Perseveration):} The most complex model, using the salience-modulated learning rule and the extended choice rule with the perseveration parameter, $\kappa$. It has four free parameters: $\alpha$, $\alpha_{\text{shift}}$, $\beta$, and $\kappa$.
\end{enumerate}

A hierarchical Bayesian approach was used, where subject-level parameters are assumed to be drawn from group-level distributions. This method allows for robust estimation of individual differences while borrowing statistical strength across the sample. All raw parameters ($\alpha_{\text{raw}}, \alpha_{\text{shift\_raw}}, \beta_{\text{raw}}, \kappa_{\text{raw}}$) were drawn from group-level normal distributions, each defined by a mean and a standard deviation. Weakly informative priors were selected for these group-level hyperparameters:
\begin{itemize}
    \item Group-level means for raw alpha and beta ($\alpha_{\mu\_\text{raw}}, \beta_{\mu\_\text{raw}}$): $\text{Uniform}(-4, 4)$
    \item Group-level mean for raw alpha-shift and kappa ($\alpha_{\text{shift}\_\mu\_\text{raw}}, \kappa_{\mu\_\text{raw}}$): $\text{Normal}(0, 1)$
    \item Group-level standard deviations for all raw parameters: $\text{Uniform}(0.001, 3)$
\end{itemize}

The models were fitted using Stan's Hamiltonian Monte Carlo sampler. Four parallel chains were run for 10,000 iterations each. The first 8,000 iterations of each chain were discarded as warmup, leaving 2,000 posterior samples per chain for analysis. Model convergence was assessed by ensuring the Gelman-Rubin statistic ($\hat{R}$) was below 1.01 for all parameters and by visual inspection of trace plots. Posterior predictive checks were planned to evaluate model fit by comparing simulated data generated from the model posteriors to the observed data.

Formal model comparison will be performed using the \texttt{loo} R package \citep{vehtari_practical_2017}. Specifically, the estimated log pointwise predictive density (elpd) will be calculated, which provides an index of a model's out-of-sample predictive accuracy using leave-one-out cross-validation. The model with the lowest elpd score (equivalent to the lowest Leave-One-Out Information Criterion, LOOIC) will be selected as the winning model. Inferences regarding the study's primary hypotheses will be based on the 95\% Highest Density Intervals (HDIs) of the posterior distributions from the winning model's parameters.

\subsubsection{Exploratory Analysis of Interindividual Differences}

To explore potential sources of interindividual variability, relationships between personality traits and model-derived parameters will be examined. The approach involves correlating the posterior mean estimates of the subject-level parameters from the winning model (e.g., $\alpha$, $\alpha_{\text{shift}}$, $\beta$, $\kappa$) with scores from the BIS-15 and SSS questionnaires (total and subscale scores). Depending on the distribution of the data, either Pearson's or Spearman's rank correlation coefficients will be calculated to probe for relationships between traits like impulsivity or sensation-seeking and individual learning or choice parameters. 

\section{Results}
\input{sections/03_results.tex}

\section{Discussion}
\input{sections/04_discussion.tex}

\printbibliography

\appendix
% Placeholder for Appendix 

% Appendix content

\section{Supplementary Materials}

\subsection{Task Instructions (English Translation)}

The following instructions were presented to participants in German. This is an English translation for reference:

\begin{quote}
Welcome to this experiment! In this task, you will make choices between two visual symbols to earn points. Your goal is to accumulate as many points as possible.

On each trial, you will see two symbols on the screen. Use the left and right arrow keys to select the symbol you want to choose. After making your choice, you will receive feedback about whether you won points or not.

Important: The probability of winning points for each symbol will change throughout the task. This means you need to pay attention to which symbol is currently better and adapt your strategy accordingly.

The position of the symbols (left or right) will change randomly, but the symbols themselves remain the same. Focus on the identity of the symbol, not its position.

Some wins will be accompanied by special audiovisual effects. These are just different ways of presenting the same win outcome.

You can earn a bonus of up to €3 based on your performance in this task. Try to earn as many points as possible!

Press SPACE to begin the practice trials.
\end{quote}

\subsection{Questionnaire Items}

\subsubsection{Barratt Impulsiveness Scale (BIS-15) - German Version}

The following items were presented with a 4-point Likert scale (1 = rarely/never, 2 = occasionally, 3 = often, 4 = almost always/always):

\begin{enumerate}
\item Ich plane Aufgaben sorgfältig. (I plan tasks carefully.)
\item Ich tue Dinge ohne nachzudenken. (I do things without thinking.)
\item Ich treffe schnelle Entscheidungen. (I make quick decisions.)
\item Ich bin unbekümmert. (I am carefree.)
\item Ich konzentriere mich nicht. (I don't concentrate.)
\item Ich habe Gedanken, die sich schnell abwechseln. (I have racing thoughts.)
\item Ich plane Reisen im Voraus. (I plan trips in advance.)
\item Ich bin selbstbeherrscht. (I am self-controlled.)
\item Ich konzentriere mich leicht. (I concentrate easily.)
\item Ich spare regelmäßig. (I save regularly.)
\item Ich winde mich und zappele bei Aufführungen oder Vorträgen. (I squirm at plays or lectures.)
\item Ich bin ein sorgfältiger Denker. (I am a careful thinker.)
\item Ich plane für Arbeitsplatzsicherheit. (I plan for job security.)
\item Ich sage Dinge ohne nachzudenken. (I say things without thinking.)
\item Ich mag es, über komplexe Probleme nachzudenken. (I like to think about complex problems.)
\end{enumerate}

\subsubsection{Sensation Seeking Scale (SSS) - German Short Version}

Participants chose between two options for each item:

\begin{enumerate}
\item 
\begin{itemize}
\item[A] Ich würde gerne Fallschirmspringen lernen.
\item[B] Ich würde niemals Fallschirmspringen wollen.
\end{itemize}

\item 
\begin{itemize}
\item[A] Ich ziehe Leute vor, die emotional ausdrucksstark sind, auch wenn sie etwas instabil sind.
\item[B] Ich ziehe Leute vor, die emotional ausgeglichen sind.
\end{itemize}

[Additional items would be listed here...]
\end{enumerate}

\subsection{Model Specifications}

\subsubsection{Stan Model Code}

The hierarchical Bayesian reinforcement learning model was implemented in Stan. The complete model code is available in the project repository at: [REPOSITORY URL]

Key model features:
\begin{itemize}
\item Hierarchical structure with group-level and subject-level parameters
\item Rescorla-Wagner learning rule with salience modulation
\item Softmax choice function
\item Weakly informative priors
\end{itemize}

\subsubsection{Prior Specifications}

Group-level priors:
\begin{itemize}
\item $\mu_\alpha \sim \text{Normal}(0, 1)$ (logit scale)
\item $\sigma_\alpha \sim \text{Half-Normal}(0, 0.5)$
\item $\mu_{\alpha_{\text{shift}}} \sim \text{Normal}(0, 0.5)$
\item $\sigma_{\alpha_{\text{shift}}} \sim \text{Half-Normal}(0, 0.5)$
\item $\mu_\beta \sim \text{Normal}(0, 1)$ (log scale)
\item $\sigma_\beta \sim \text{Half-Normal}(0, 0.5)$
\end{itemize}

\subsection{Additional Analyses}

\subsubsection{Model Comparison}

Alternative models were fitted to assess the necessity of the salience modulation parameter:

\begin{enumerate}
\item \textbf{Null Model:} Standard RL model without salience effects
\item \textbf{Full Model:} RL model with salience-modulated learning rate
\item \textbf{Alternative Model:} RL model with salience affecting choice temperature
\end{enumerate}

Model comparison was conducted using leave-one-out cross-validation (LOO-CV) and the Widely Applicable Information Criterion (WAIC).

\subsubsection{Sensitivity Analyses}

Robustness of findings was assessed through:
\begin{itemize}
\item Alternative prior specifications
\item Exclusion of participants with high miss rates (>10\%)
\item Analysis of different trial subsets (early vs. late trials)
\end{itemize}

\subsection{Data and Code Availability}

\subsubsection{Data Sharing}

Anonymized data supporting the conclusions of this study are available upon reasonable request to the corresponding author, subject to ethical approval and data protection regulations.

\subsubsection{Code Repository}

All analysis code, including data preprocessing, model fitting, and visualization scripts, is available at: [REPOSITORY URL]

The repository includes:
\begin{itemize}
\item R scripts for data analysis
\item Stan model files
\item Python scripts for data collection
\item Documentation and README files
\end{itemize}

\subsection{Ethical Approval Documentation}

This study was approved by the Ethics Committee of the Faculty of Human Sciences, University of Cologne (Approval No. JPHF0271, granted December 11, 2024). The study was conducted in accordance with the Declaration of Helsinki and local institutional guidelines.

\subsection{Preregistration}

This study was preregistered prior to data collection. The preregistration document is available at: [PREREGISTRATION URL]

\section{Declaration of Independent Work}

I hereby declare that this master's thesis has been written independently and without unauthorized assistance. All sources and aids used have been properly cited and acknowledged. This work has not been submitted for any other academic degree or professional qualification.

\vspace{2cm}

\noindent
Cologne, \today

\vspace{2cm}

\noindent
\rule{6cm}{0.4pt}\\
Yannik Poth 

\end{document} 